Un radiateur à eau chaude a une aire de surface de \qty{\v{S}}{m^2} et est recouvert d'une peinture d'émissivité \num{\v{em}}. Sa surface est à \qty{\v{Tr}}{\degreeCelsius} et les murs de la pièce sont à \qty{\v{Te}}{\degreeCelsius}. On ne s'intéresse qu'aux échanges par rayonnement.

\begin{enumerate}[label=\alph*)]
    \item Déterminez la puissance nette que le radiateur fournit à la pièce par rayonnement.
    \item À quelle température faudrait-il porter le radiateur pour qu'il fournisse \qty{\v{Pv}}{W} nets par rayonnement ?
    \item Pourquoi une peinture métallisée (aluminium, émissivité faible) est-elle un mauvais choix pour un radiateur ?
\end{enumerate}
